\setcounter{aufgabe}{0}
\hypertarget{e1}{}\einheitenkopf{Einheit 1 von 5 · Proportionale Funktion}
\zweigzeile{Hier lernst du, Funktionen mit $y = m \cdot x$ zu berechnen, zu zeichnen und zu erkennen · neu in diesem Jahr · P10 oft · baut auf: Koordinatensystem, proportionale Zuordnungen (Dreisatz)}

\begin{aufgabe}{\textbf{Ich kann die Gerade zu $y = m \cdot x$ zeichnen.} Berechne $y$ für $x = 1$. Trage den Punkt ein und zeichne die Gerade durch den Ursprung $O(0\,|\,0)$ und diesen Punkt. Zeichne a) und b) links, c) und d) rechts.}
\beispiel{y = 2x:\quad y &= 2 \cdot 1 = 2 && \text{Punkt } P(1\,|\,2),\ \text{Gerade durch } O \text{ und } P}
\begin{teilezwei}
\tz $y = x$ \quad $x = 1$: \feld{y} & \tz $y = 3x$ \quad $x = 1$: \feld{y} \\
\tz $y = 4x$ \quad $x = 1$: \feld{y} & \tz $y = 5x$ \quad $x = 1$: \feld{y} \\
\end{teilezwei}

\begin{ksys}[xmin=-1,xmax=3,ymin=-1,ymax=6]
\end{ksys}\ksysabstand\begin{ksys}[xmin=-1,xmax=3,ymin=-1,ymax=6]
\end{ksys}

\end{aufgabe}

\begin{aufgabe}{\textbf{Ich kann die Gerade zu $y = m \cdot x$ zeichnen – weiter.} Berechne einen Punkt, trage ihn ein und zeichne die Gerade durch den Ursprung. Zeichne alle drei Geraden in dasselbe Koordinatensystem.}
\setcounter{teil}{4}
\begin{teile}
\teil $y = 2{,}5x$ \quad $x = 2$: \feld{y}
\teil $y = \frac{1}{3}x$ \quad Wähle $x$ so, dass $y$ eine ganze Zahl ist: \quad \feld{x} \feld{y}
\teil $y = -2x$ \quad $x = 1$: \feld{y}
\end{teile}

\begin{ksys}[xmin=-2,xmax=4,ymin=-4,ymax=6]
\end{ksys}

\end{aufgabe}

\begin{aufgabe}{\textbf{Ich kann Werte an einer Ursprungsgeraden ablesen und berechnen.} Der Graph gehört zur Funktion $y = 1{,}5 \cdot x$.}

\begin{ksys}[xmin=-3,xmax=5,ymin=-4,ymax=7,ablesen]
\gerade[4.5]{1.5}{0}{f}
\end{ksys}

Lies am Graphen ab.
\begin{teile}
\teil Welcher $y$-Wert gehört zu $x = 2$? \quad \feld{y}
\teil Welcher $y$-Wert gehört zu $x = -2$? \quad \feld{y}
\teil Welcher $x$-Wert gehört zu $y = 6$? \quad \feld{x}
\end{teile}
Rechne mit der Gleichung.
\begin{teile}
\teil Rechne a) mit der Gleichung nach. Stimmt dein abgelesener Wert?
\teil Berechne $y$ für $x = 12$. Dieser Punkt liegt nicht mehr im Bild. \quad \feld{y}
\teil Prüfe e) mit dem Dreisatz. Beginne mit dem Wertepaar aus a).
\end{teile}
\begin{dreisatz}{$x$}{$y$}
\dsz{2}{3}
\dsp{\cdot 6}{\cdot 6}
\dsz{12}{\dsleer}
\end{dreisatz}
\end{aufgabe}

\begin{aufgabe}{\textbf{Ich kann erkennen, ob eine proportionale Funktion vorliegt, und ihre Gleichung angeben.} Entscheide, ob $y$ proportional zu $x$ ist. Wenn ja, gib die Gleichung an.}
\begin{teile}
\teil \wertetabelle[6,12,18,24]{x}{y}{1,2,3,4} \par \janein \quad \feldl{y}
\teil \wertetabelle[7,14,20]{x}{y}{2,4,6} \par \janein \quad \feldl{y}
\teil \wertetabelle[5,8,11]{x}{y}{1,2,3} \par \janein \quad \feldl{y}
\teil \wertetabelle[3,-1{,}5,-6]{x}{y}{-2,1,4} \par \janein \quad \feldl{y}
\teil Ein Liter Saft kostet 1,80 €. $x$ ist die Anzahl der Liter, $y$ der Preis in €. \\ \janein \quad \feldl{y}
\teil Eine Taxifahrt kostet 4 € Grundgebühr und 2 € je Kilometer. $x$ ist die Strecke in km, $y$ der Preis in €. \\ \janein \quad \feldl{y}
\end{teile}
\end{aufgabe}

\begin{aufgabe}{\textbf{Ich finde den Fehler beim Bestimmen des Faktors.} Tom soll zur Tabelle den Faktor $m$ und die Gleichung bestimmen.}

\wertetabelle[21,42]{x}{y}{3,6}

Tom rechnet so:
\rechnung{m &= 21 - 3 = 18 \\ y &= 18x}
\begin{teile}
\teil Finde den Fehler, benenne ihn und korrigiere die Rechnung.
\teil Bestimme den Faktor $m$ und die Gleichung zu dieser Tabelle.
\end{teile}

\wertetabelle[18,45]{x}{y}{4,10}
\end{aufgabe}

\begin{aufgabe}{\textbf{Ich kann Eigenschaften von Ursprungsgeraden begründen.} Begründe ohne Rechnung.}
\begin{teile}
\teil Warum geht der Graph jeder Funktion $y = m \cdot x$ durch den Ursprung $O(0\,|\,0)$?
\teil Kann der Punkt $P(2\,|\,5)$ auf der Geraden $y = -4x$ liegen? Entscheide und begründe.
\end{teile}
\end{aufgabe}

\begin{aufgabe}{\textbf{Ich kann eine Sachsituation mit $y = m \cdot x$ beschreiben.} In einen leeren Pool fließen gleichmäßig 40 Liter Wasser pro Minute. $x$ ist die Zeit in Minuten, $y$ die Wassermenge in Litern.}
\begin{teile}
\teil Stelle die Gleichung auf. \quad \feldl{y}
\teil Zeichne den Graphen für die ersten 10 Minuten.
\teil Was bedeutet die Zahl 40 im Sachzusammenhang? Lies am Graphen ab, nach wie vielen Minuten 200 Liter im Pool sind.
\end{teile}

\begin{ksys}[xmin=0,xmax=10,ymin=0,ymax=400,xstep=1,ystep=50,xlabel=$x$ in min,ylabel=$y$ in l]
\end{ksys}

\end{aufgabe}
